\section{Resultados y discusión}\label{sec:resultados}

\textbf{Error en validación cruzada.} El modelo obtiene un RMSE de 0.458\,t/ha, un error absoluto medio (MAE) de
0.312\,t/ha y un coeficiente de determinación ($R^2$) de 0.71~\cite{doc_modelo}. Solo con la media de cada régimen el
RMSE es 0.489. La mejora va en la dirección correcta, aunque con 138 parcelas su intervalo de confianza del 95\,\%
todavía incluye el cero (anexo técnico, Tabla~\ref{A-tab:error}).

\textbf{Error esperado en la calificación.} Al repetir toda la elección dentro de la validación el RMSE es 0.490,
prácticamente el de la media por régimen. Por eso esperamos alrededor de 0.49\,t/ha en las 59 parcelas, y el 0.458 es
una cifra optimista. Con solo 59 parcelas el resultado también depende de cuáles toquen, y por azar podría caer entre
0.35 y 0.56~\cite{doc_modelo}.

\textbf{Qué aportó cada pieza.} El régimen es la base. Cuando la validación deja fuera zonas completas, la media por
municipio sube a 0.955\,t/ha y la del régimen se mantiene en 0.489, así que el régimen generaliza y el municipio no. El
rezago espacial mejoró el error en 9 de 10 repeticiones y pasó la prueba de permutación ($p=0.02$). Ningún bloque de
variables satelitales, climáticas o topográficas pasó la regla (Fig.~\ref{A-fig:bloques}). El clima, por ejemplo,
viene en celdas de unos 5\,km y solo distingue zonas, algo que ya hacen el régimen y los vecinos. Entre los algoritmos
comparados, el proceso gaussiano obtuvo el menor error~\cite{doc_modelo}.

\textbf{Cómo verificamos la fuga de información.} Auditamos que todo lo que usa el rendimiento se calcule solo con el
entrenamiento, una prueba automatizada verifica que cambiar las etiquetas de prueba no cambia sus predicciones, y con
los rendimientos revueltos el RMSE no baja de 0.485, contra 0.467 con los reales~\cite{doc_modelo}. Ninguna prueba
descarta la fuga por completo, pero ninguna la encontró.

\textbf{Los intervalos cumplen.} El intervalo del percentil 10 al 90 contiene el rendimiento real en 83\,\% de los
casos, cerca de la meta de 80\,\%~\cite{doc_modelo}. En 18 parcelas, todas del régimen B, la probabilidad de
quedar bajo 2.5\,t/ha, el percentil 10 de ese régimen, supera 10\,\%. La tabla con las 59 predicciones está en el
anexo técnico.

\textbf{El modelo municipal funciona donde no hay datos.} En municipios sin datos del SIAP a menos de 50\,km obtiene un
RMSE de 0.44\,t/ha en promedios municipales, que no se compara con el de parcelas, contra 0.57 del kriging de vecinos,
que no usa clima, y gana en las 10 repeticiones (anexo técnico, Fig.~\ref{A-fig:municipal_modelos}). Como elegimos la
Tmax viendo resultados, repetimos esa elección de forma anidada y la ventaja sigue a favor, con unos 0.07\,t/ha. Cada
grado más de Tmax se asocia con 0.79\,t/ha menos~\cite{doc_generalizable}. El mapa
regional (Fig.~\ref{A-fig:mapas}) muestra que el oriente de Puebla y Tlaxcala, más fresco en verano, tiene el mayor
rendimiento esperado.

\textbf{Por qué el clima no ayudó en las parcelas.} En los nueve municipios del reto la Tmax solo va de 23.3 a
23.9\,$^\circ$C~\cite{doc_generalizable}. A esa escala casi no varía, así que no puede distinguir una parcela de otra.
A escala regional sí.

\textbf{Preguntas guía.} Las bases plantean tres preguntas~\cite{bases}. Como el modelo final no usa índices ni
clima, las respondimos con un bosque aleatorio auxiliar, midiendo cuánto sube su error al revolver cada variable
(anexo técnico, Fig.~\ref{A-fig:indices}). Primero, el índice satelital más importante. Sin el régimen es el NDVI,
porque su historia identifica la zona, y con régimen y vecinos es el índice de sequía NDDI~\cite{gu2007}, aunque apenas
aporta 0.004\,t/ha~\cite{doc_modelo}. Segundo, la variable climática de mayor peso. En las parcelas ninguna aporta una
vez conocido el régimen, y a escala regional es la Tmax de julio y agosto. Tercero, la mejor combinación, que es
régimen más rezago espacial.

\textbf{Qué significa el error para un productor.} Con el precio del SIAP, un error de 0.49\,t/ha equivale a unos
2,700 pesos por hectárea en promedio por parcela, entre 2,100 y 3,500 según el municipio. Es de 7 a 13\,\% del costo de
producción y de 14 a 47\,\% del crédito por hectárea~\cite{repo}. Sirve para ordenar parcelas por riesgo, pero no para
decidir sola sobre una parcela.

\textbf{Aplicación.} En \url{https://agrocebada.alan-pv.com} se consulta el mapa de parcelas, la predicción con su
intervalo y la probabilidad de quedar bajo un umbral que elige el usuario. La aplicación ejecuta el modelo en vivo.
